%% world_model_paper.tex
%% Dependently Typed World Models: Toward Proof-Carrying Learned Simulators
%% Draft — August 2026

\documentclass[12pt, a4paper]{article}

\usepackage[T1]{fontenc}
\usepackage[utf8]{inputenc}
\usepackage[margin=1.25in]{geometry}
\usepackage{amsmath, amssymb, amsthm}
\usepackage{mathtools}
\usepackage[expansion=false]{microtype}
\usepackage{parskip}
\usepackage[hidelinks]{hyperref}
\usepackage{listings}
\usepackage{xcolor}
\usepackage{booktabs}
\usepackage{caption}
\usepackage{natbib}
\usepackage{abstract}
\usepackage{titlesec}

% ---- typography ----
\titleformat{\section}{\large\bfseries}{\thesection.}{0.6em}{}
\titleformat{\subsection}{\normalsize\bfseries}{\thesubsection.}{0.5em}{}
\setlength{\parindent}{1.4em}
\setlength{\parskip}{0pt}

% ---- listings style for Rocq / Coq ----
\definecolor{rocqblue}{RGB}{0,80,160}
\definecolor{rocqgray}{RGB}{100,100,100}
\definecolor{rocqgreen}{RGB}{0,120,60}
\definecolor{panelbg}{RGB}{236,244,234}

\lstdefinelanguage{Rocq}{
  keywords={From,Require,Import,Section,Context,Record,Definition,Lemma,Proof,
            Qed,End,forall,exists,fun,match,with,let,in,if,then,else,
            Inductive,Fixpoint,Variable,Hypothesis,Axiom,Notation,
            move,rewrite,case,apply,exact,intro,intros,simpl,auto,by},
  keywordstyle=\color{rocqblue}\bfseries,
  comment=[l]{(*},
  commentstyle=\color{rocqgray}\itshape,
  stringstyle=\color{rocqgreen},
  morestring=[b]",
  sensitive=true,
  basicstyle=\ttfamily\small,
  backgroundcolor=\color{panelbg},
  frame=single,
  framerule=0.4pt,
  rulecolor=\color{rocqgray!60},
  breaklines=true,
  captionpos=b,
  xleftmargin=1.2em,
  xrightmargin=0.6em,
  aboveskip=0.8em,
  belowskip=0.6em
}

\lstset{language=Rocq}

% ---- theorem environments ----
\newtheorem{definition}{Definition}[section]
\newtheorem{proposition}{Proposition}[section]
\newtheorem{remark}{Remark}[section]

% ---- metadata ----
\title{\textbf{Dependently Typed World Models}\\[0.4em]
  \large Toward Proof-Carrying Learned Simulators}
\author{Jordan Rule\\[0.2em]
  \small \texttt{jrule@uchicago.edu}}
\date{August 2026}

\begin{document}

\maketitle
\thispagestyle{empty}

% ---- abstract ----
\begin{abstract}
\noindent
World models are the internal working pictures through which modern learning
systems rehearse possible futures. They compress experience, imagine
consequences, and help an agent choose before it acts. Yet their most
important commitments---what a latent state is allowed to mean, when evidence
is sufficient for planning, and which actions are legitimate under uncertainty
---are often left implicit. This paper argues that dependent type theory offers
a practical way to make those commitments explicit and machine-checkable. We
introduce the notion of a \emph{proof-carrying world model}: a learned
simulator whose states, predictions, and action proposals carry evidence that
their preconditions are met. We connect this view to Joint Embedding Predictive
Architecture (JEPA), where learning already proceeds from partial views toward
predictive sufficiency rather than exhaustive reconstruction. The result is a
research program in which world models remain empirical and data-driven while
becoming clearer about admissibility, provenance, and safe use in multi-agent
settings.
\end{abstract}

\bigskip
\noindent\textit{Keywords:} world models, dependent type theory, Rocq, JEPA,
model-based reinforcement learning, proof-carrying programs, alignment,
multi-agent systems.

\newpage

%% ─────────────────────────────────────────────────────────────────────────────
\section{Introduction}
%% ─────────────────────────────────────────────────────────────────────────────

One of the most consequential ideas in artificial intelligence is that an agent
need not meet the world only in the instant. It may, before acting, carry an
inner stage on which candidate futures can be rehearsed. In model-based
reinforcement learning, that stage is the world model: a learned internal
simulator that compresses experience into a latent state and uses that state to
forecast what actions might bring about
\citep{sutton1990dyna, ha2018world, hafner2019dream, schrittwieser2020mastering,
hafner2023mastering, ye2021efficientzero}. The practical appeal is easy to see.
An agent that can imagine cheaply does not need to experiment as recklessly in
reality.

The idea has also widened. In one lineage, world models are used explicitly for
planning. In another, adjacent lineage, predictive systems such as Joint
Embedding Predictive Architecture (JEPA) learn latent representations by asking
what one partial view of a situation should allow a system to say about another
\citep{lecun2022path, assran2023ijepa}. These models do not insist on
reconstructing every pixel or every sensory detail. They aim instead for a more
selective kind of understanding: enough internal structure to support
prediction, abstraction, and downstream decision-making.

And yet the central mystery remains curiously underdescribed. A latent state in
most current systems is useful, trainable, and often powerful, but semantically
reticent. It carries little machine-checkable explanation of what it is allowed
to represent, which fragments of evidence may be combined into it, or what
constraints its imagined futures are supposed to respect. These commitments, if
they are written down at all, typically live outside the model—in design docs,
policy layers, runtime checks, or after-the-fact audits.

This paper argues that dependent type theory offers a way to move some of that
structure inward. In a dependently typed setting, a value can travel together
with evidence that it satisfies the conditions required for its use. Proof need
not arrive as a bureaucratic seal applied later; it can accompany the object at
the moment of construction. We draw in particular on the Rocq proof assistant
\citep{coq2024}, using dependent records and constructive proofs to make world
model interfaces explicit.

Before any formal machinery, however, there is a gentler picture. One can think
of a world model as something assembled from partial views: a few observations
here, a capability hint there, a short trace of what happened when an action was
taken. Some views fit together; some contradict one another; some are simply too
thin to justify high-confidence plans. This patchwork intuition is simple enough
to grasp in ordinary language, yet it already echoes JEPA's core move: learn to
predict what is hidden from what is seen, without pretending to know everything.
Our claim is that dependent types can give this predictive practice a clearer
grammar of admissible use.

Our contributions are:
\begin{enumerate}
  \item A conceptual framework---the \emph{proof-carrying world model}---that
    articulates how dependent types can annotate the operational interfaces of
    a learned simulator (\S\ref{sec:framework}).
  \item A practical bridge from partial views to typed predictive contracts,
    showing how JEPA-style representation learning can be coupled to explicit
    admissibility obligations for planning under uncertainty
    (\S\ref{sec:bridge}).
  \item Small Rocq formalizations illustrating these ideas, from basic
    proof-carrying state through typed context-target prediction and
    evidence-aware action proposals (\S\ref{sec:rocq}).
  \item A discussion of research directions, including explicit connections to
    alignment-aware planning and multi-agent trust (\S\ref{sec:discussion}).
\end{enumerate}

We do not claim a complete implementation. We offer a research direction and
a working vocabulary for future development.

Figure~\ref{fig:agenda-stack} provides a compact visual of the agenda: local
predictive artifacts feed shared admissibility checks, which in turn stabilize
a JEPA-style planning interface.

\begin{figure}[t]
\centering
\fbox{%
\begin{minipage}{0.92\linewidth}
\centering
\small
\textbf{Research Agenda Stack: From Local Predictive Artifacts to a Verified JEPA Interface}\\[0.6em]
\begin{tabular}{c}
\fbox{\parbox{0.84\linewidth}{\centering
Distributed teams publish local artifacts\\
$\mathcal{A}_i = (\text{partial views}_i,\, \text{typed obligations}_i,\, \text{predictive patches}_i)$
}} \\
$\Downarrow$ \\
\fbox{\parbox{0.84\linewidth}{\centering
Shared Rocq refinement checks align admissibility rules across artifacts
}} \\
$\Downarrow$ \\
\fbox{\parbox{0.84\linewidth}{\centering
Joint embedding training uses only contract-satisfying context-target pairs
}} \\
$\Downarrow$ \\
\fbox{\parbox{0.84\linewidth}{\centering
Formally constrained JEPA interface for planning and rollout
}}
\end{tabular}
\end{minipage}%
}
\caption{Illustrative pipeline for a globally collaborative program. The claim is
not that learning disappears into logic, but that learned latent updates are
filtered through shared, explicit admissibility contracts.}
\label{fig:agenda-stack}
\end{figure}

%% ─────────────────────────────────────────────────────────────────────────────
\section{Background}
\label{sec:background}
%% ─────────────────────────────────────────────────────────────────────────────

\subsection{World Models in Reinforcement Learning}

A world model, in the narrow reinforcement-learning sense, is a learned
surrogate for the dynamics of an environment. Formally one may write it as a
map $f_\theta : \mathcal{Z} \times \mathcal{A} \to \mathcal{Z} \times
\mathcal{R}$, sending a latent state and an action to a predicted successor
state and reward. But that notation only hints at the broader intuition. A
world model is a compressed, internal picture of how things change: what tends
to follow what, which possibilities remain open, and which imagined futures are
worth spending computation on.

The modern lineage is familiar. Dyna first made explicit the advantage of
learning and planning in tandem \citep{sutton1990dyna}. Ha and Schmidhuber's
World Models popularized the image of an agent dreaming inside its own latent
space \citep{ha2018world}. Dreamer and its successors showed that long-horizon
imagination could become a practical engine for control
\citep{hafner2019dream, hafner2021mastering, hafner2023mastering}. MuZero and
EfficientZero demonstrated that one need not reconstruct the visible world in
full detail in order to plan effectively \citep{schrittwieser2020mastering,
ye2021efficientzero}.

JEPA extends the same family resemblance in a slightly different direction.
Rather than learning to regenerate an entire input, a JEPA-style system learns
to predict the representation of a withheld target view from an available
context view \citep{lecun2022path, assran2023ijepa}. What matters is not visual
fidelity but predictive sufficiency: the latent state should preserve the
structure needed to say something reliable about what lies beyond the current
glimpse. In that sense, JEPA is not a departure from the world-model story so
much as a sharpening of one of its oldest instincts.

Across these traditions, a recurring tension remains between \emph{predictive
skill} and \emph{explicit meaning}. A model may be extraordinarily good at
forecasting what happens next and still be silent about what sort of state it
is claiming to inhabit, what sources of evidence it has mixed together, and
what operational limits should govern its use. This gap is the starting point
for our proposal.

\subsection{Dependent Type Theory and the Rocq Proof Assistant}

Dependent type theory \citep{martin-lof1984, coquand1988calculus} begins from a
simple but powerful thought: the description of an object may be allowed to
mention the object itself. Instead of saying merely ``this is a vector,'' one
may say ``this is a vector of length $n$,'' and insist that $n$ be part of the
object's type. A familiar example is $\mathsf{Vec}(n)$, the type of vectors of
exactly $n$ elements. More generally, a value of type $\{x : A \mid P(x)\}$ is
both a value in $A$ and evidence that it satisfies $P(x)$.

For readers outside formal methods, the philosophical point matters more than
the notation. In an ordinary program, one often constructs a value first and
checks later whether it meets the relevant conditions. In a dependently typed
program, one can require those conditions at birth. Some classes of mismatch
never become runtime problems because the system refuses to construct the wrong
object in the first place.

Rocq (formerly Coq) \citep{coq2024} implements the Calculus of Inductive
Constructions \citep{coquand1988calculus, paulin-mohring2015}, a rich
dependent type theory supporting:
\begin{itemize}
  \item \emph{Inductive types} and pattern matching;
  \item \emph{Dependent records} (\texttt{Record}) that bundle values with
    their correctness conditions;
  \item \emph{Proof terms} as first-class citizens of the same language as
    programs;
  \item \emph{Extraction} to OCaml or Haskell, enabling certified programs to
    run outside the proof environment.
\end{itemize}

The MathComp library \citep{gonthier2013odd} provides mature foundations for
constructive formalization in Rocq and practical proof engineering at scale.

\subsection{Partial Views, Predictive Sufficiency, and Contracts}

Before introducing heavy formal structure, it helps to start with a familiar
scene. An agent rarely observes a complete world. It sees fragments: a camera
crop, a noisy sensor packet, a short action history, a partial map of who said
what. JEPA-style learning asks what one such fragment should allow us to infer
about another \citep{lecun2022path, assran2023ijepa}. This is a remarkably
practical stance for autonomous systems: do not reconstruct everything; recover
enough structure to make the next decision responsibly.

Dependent types sharpen that stance. They let us say: this prediction is usable
for planning only when certain evidence conditions hold. Instead of treating
latent vectors as universally admissible, we can treat them as claims with scope.
Some claims are broad enough for low-risk suggestion. Others are narrow and
should not authorize consequential actions. The technical language is precise,
but the intuition is plain: uncertainty is not an afterthought; it is part of
the object we manipulate.

\subsection{From JEPA Representations to Typed Planning}

In many deployments, the central failure mode is not poor prediction in isolation
but overconfident use of prediction in the wrong context. A model can be right on
average and still dangerous in the moments that matter. The bridge we propose is
therefore interface-level: keep JEPA as the statistical engine for partial-view
prediction, and place a typed contract at the boundary where predictions become
plans.

Under this view, a planner does not merely consume a latent state. It consumes a
\emph{checked} latent state, paired with explicit witnesses about provenance,
coverage, or confidence regime. The point is not to freeze learning into rigid
symbolism; it is to stop treating all latent states as equally action-worthy.

% ---- formal development ----
\section{A Framework for Proof-Carrying World Models}
\label{sec:framework}
%% ─────────────────────────────────────────────────────────────────────────────

\subsection{Motivation}

Most learned latent states are described, at bottom, as arrays of numbers. That
description is enough for optimization, but it says very little about meaning.
It does not say whether the state was assembled from compatible evidence,
whether it presupposes a certain capability scope, or whether it is even the
kind of state a planner should be allowed to imagine from. Those judgments are
usually external to the model.

Our proposal is to make admissibility \emph{intrinsic}. If a state is meant to
stand for ``a history built from non-conflicting observations under capability
scope $\kappa$,'' then its type should say so. If a plan is meant to require
certain permissions, the need for those permissions should be visible in the
plan's construction. Scope confusion then ceases to be a policy bug discovered
late; it becomes a type error discovered early.

\subsection{Proof-Carrying State}

We define a \emph{proof-carrying state} as a dependent record bundling a
representational value with a machine-checked admissibility witness.

\begin{definition}[Proof-Carrying State]
Let $S$ be a set of representational states and let
$\mathsf{Adm} : S \to \mathsf{Prop}$ be an admissibility predicate. A
\emph{proof-carrying state} is a pair $(s, h)$ where $s \in S$ and $h :
\mathsf{Adm}(s)$.
\end{definition}

The key property is psychological as well as technical. One is no longer asked
to trust that a state has been blessed somewhere upstream. The blessing is part
of the object. Construction and certification occur together, or not at all.

\subsection{Proof-Carrying World Model Interface}

A world model operating over proof-carrying states must be designed so that:
\begin{enumerate}
  \item \emph{Observation closure}: applying the observation operator to an
    admissible state produces an admissible state.
  \item \emph{Imagination validity}: applying the imagination (rollout)
    operator to an admissible state produces a valid (non-undefined) state.
\end{enumerate}

These are not runtime checks but proof obligations discharged at definition
time. A world model that cannot satisfy them is, in a precise sense,
\emph{ill-typed}—it cannot be instantiated.

The following Rocq sketch captures the interface:

\begin{lstlisting}[caption={Proof-carrying world model interface in Rocq},
                   label={lst:world-interface}]
From mathcomp Require Import ssreflect ssrbool ssrnat seq.

Section TypedWorld.
Context {State Observation Action : Type}.
Context (Adm : State -> Prop).

Record CheckedState := {
  state_val : State;
  state_ok  : Adm state_val
}.

Record WorldModel := {
  observe : CheckedState -> seq Observation -> CheckedState;
  imagine : CheckedState -> seq Action -> option CheckedState
}.

End TypedWorld.
\end{lstlisting}

This interface is intentionally schematic. The key point is that planning-time
outputs are partial (\texttt{option}) and must justify admissibility on success.
A JEPA-style model can be read through the same lens: context builds a latent
state, prediction proposes a target representation, and typed obligations specify
when that representation is safe to use for downstream planning.

%% ─────────────────────────────────────────────────────────────────────────────
\section{From Partial Views to Planning Contracts}
\label{sec:bridge}
%% ─────────────────────────────────────────────────────────────────────────────

\subsection{JEPA and the Grammar of Missing Information}

JEPA's conceptual strength is that it takes missing information seriously. It
asks what can be inferred from what is present, without pretending that hidden
structure has been directly observed \citep{lecun2022path, assran2023ijepa}.
For autonomous systems operating in the wild, this is the norm rather than the
exception: data arrive late, sensors disagree, and collaborators report through
uneven trust channels.

Dependent types offer a useful complement: they let us distinguish between
``predicted'' and ``safely actionable'' in the type of an object, not only in a
policy note outside the model. This creates a clean handoff: JEPA supplies
statistical abstraction; typed contracts govern admissible use.

\subsection{Planning Under Limited Knowledge}

When knowledge is partial, planning should degrade gracefully rather than fail
silently or overreach. A typed interface can enforce this behavior by requiring
that high-impact actions carry explicit witnesses that their preconditions are
met under the current evidence profile. If those witnesses are absent, the plan
construction itself fails.

This is not merely elegant formalism. It is operationally relevant in settings
where autonomous systems must choose under ambiguity. A planner that cannot prove
enough should return an abstention, a request for more context, or a lower-risk
fallback policy. Typed construction gives these behaviors first-class status.

\subsection{A Collaborative Research Loop}

Figure~\ref{fig:global-jepa-loop} sketches the intended workflow. Distributed
research groups contribute partial-view datasets, local obligations, and
refinement lemmas. Shared checks keep only artifacts that satisfy the agreed
contracts. JEPA-style training then updates representations on this filtered
corpus, and the next iteration tightens both statistical quality and formal
clarity.

\begin{figure}[t]
\centering
\fbox{%
\begin{minipage}{0.92\linewidth}
\centering
\small
\textbf{Global Collaboration Loop: Typed Publication to JEPA Update}\\[0.6em]
\begin{tabular}{c}
Iteration $k$ \\
\fbox{\parbox{0.84\linewidth}{\centering
Teams publish $(\text{partial views},\, \text{local proofs},\, \text{context-target pairs})_k$
}} \\
$\Downarrow$ \\
\fbox{\parbox{0.84\linewidth}{\centering
Common admissibility checker: keep only contract-satisfying artifacts
}} \\
$\Downarrow$ \\
\fbox{\parbox{0.84\linewidth}{\centering
JEPA update on filtered corpus: $E_k \mapsto E_{k+1}$
}} \\
$\Downarrow$ \\
\fbox{\parbox{0.84\linewidth}{\centering
Publish stronger obligations and refinement lemmas for iteration $k+1$
}}
\end{tabular}
\end{minipage}%
}
\caption{Indicative training-governance cycle. Representation quality improves
statistically while admissibility guarantees improve constructively.}
\label{fig:global-jepa-loop}
\end{figure}

\subsection{Immediate Alignment Value in Multi-Agent Settings}

The safety advantage of this bridge is concrete. In multi-agent environments,
failures often arise not from a single bad prediction but from a mismatch between
what one agent inferred and what another agent was entitled to do. Typed world
model interfaces make these mismatches visible at construction time. Provenance,
permission scope, and uncertainty regime can all be made part of the data that
planning consumes.

The resulting architecture remains learned and adaptive, but it becomes less
willing to improvise beyond its evidence. That is precisely the behavior we want
from autonomous systems that must act among other autonomous systems.

%% ─────────────────────────────────────────────────────────────────────────────
\section{Illustrative Rocq Formalizations}
\label{sec:rocq}
%% ─────────────────────────────────────────────────────────────────────────────

\subsection{A Tiny Example of Proof-Carrying State}

Before introducing additional interface structure, we illustrate the simplest instance of
proof-carrying state: a rollout record whose \texttt{horizon} field is
certified by a length witness.

\begin{lstlisting}[caption={Minimal proof-carrying rollout record},
                   label={lst:minimal}]
From Coq Require Import List Arith.
Import ListNotations.

Record CheckedRollout := {
  horizon         : nat;
  actions         : list nat;
  horizon_matches : length actions = horizon
}.

Definition single_action (a : nat) : CheckedRollout :=
  {| horizon         := 1;
     actions         := [a];
     horizon_matches := eq_refl |}.
\end{lstlisting}

The field \texttt{horizon\_matches} is not a runtime assertion but a proof
term. It is discharged by \texttt{eq\_refl} at the point of construction,
confirming that the singleton list \texttt{[a]} has length~1. A
\texttt{CheckedRollout} with mismatched horizon cannot be constructed; the
attempt would produce a type error. This is the paper's basic promise in
miniature: if a state claims to summarize a horizon, that claim can be made
part of the object rather than an informal comment about it.

\subsection{Typed Context-Target Prediction}

The following Rocq definition illustrates a proof-carrying context-target
prediction pair. The key point is that the prediction is only usable when the
context is known to be non-empty.

\begin{lstlisting}[caption={Typed context-target prediction record},
                   label={lst:rollout}]
From Coq Require Import List Arith.
Import ListNotations.

Record PredictivePair := {
  context_tokens : list nat;
  target_tokens  : list nat;
  context_nonempty : length context_tokens > 0
}.

Definition safe_predict (p : PredictivePair) : option (list nat) :=
  if Nat.eqb (length (context_tokens p)) 0
  then None
  else Some (target_tokens p).
\end{lstlisting}

The \texttt{safe\_predict} function is deliberately modest: prediction is allowed
only when the context witness says there is actually context to condition on. The
point is not sophistication but shape. Even this tiny example reflects the design
rule that planning-time prediction should expose admissibility directly.

\subsection{Evidence-Scoped Action Proposals}

A second pattern is to type action proposals by the evidence they require. Rather
than ``choose action,'' one defines ``choose action under evidence profile $e$.''
When the profile is weak, the system can still reason, but only into action
families allowed for weak evidence.

\begin{lstlisting}[caption={Evidence-scoped action proposal},
                   label={lst:sensor}]
From Coq Require Import List Arith.

Inductive RiskLevel := Low | High.

Record Evidence := {
  confidence : nat;
  confidence_ok : confidence <= 100
}.

Definition can_take_high_risk (e : Evidence) : bool :=
  Nat.leb 80 (confidence e).

Definition propose_action (e : Evidence) : RiskLevel :=
  if can_take_high_risk e then High else Low.
\end{lstlisting}

The program is simple, but the design lesson is durable: evidence quality becomes
an explicit input to action construction. In richer systems, those thresholds are
replaced by proofs over learned uncertainty estimates and provenance contracts.

\subsection{Composing Multi-Agent Signals with Explicit Trust}

The final pattern is trust-aware composition. In cooperative and adversarial
settings alike, a planner must know not only \emph{what} was observed, but
\emph{who observed it} and under what trust assumptions.

\begin{lstlisting}[caption={Trust-annotated observation packet},
                   label={lst:compose}]
From Coq Require Import List Arith.

Inductive Trust := Unverified | Verified.

Record Observation := {
  payload : nat;
  trust   : Trust
}.

Definition usable_for_critical_plan (o : Observation) : bool :=
  match trust o with
  | Verified => true
  | Unverified => false
  end.
\end{lstlisting}

This sketch captures the core operational idea: critical planning paths should be
typed to require stronger evidence than exploratory paths. The architecture can
still learn from all signals, but it need not treat all signals as equally
authoritative at decision time.

%% ─────────────────────────────────────────────────────────────────────────────
\section{Discussion and Research Directions}
\label{sec:discussion}
%% ─────────────────────────────────────────────────────────────────────────────

\subsection{Alignment-Aware Planning}

The alignment problem for learned agents \citep{russell2019human,
gabriel2020alignment} is, in part, a problem of specification: how does an
agent know that the actions it is considering are within the scope of what it
is supposed to do? Current approaches rely heavily on reward shaping, RLHF
\citep{christiano2017rlhf, bai2022training}, and constitutional constraints
\citep{bai2022constitutional}, all of which are intrinsically extrinsic—they
surround the model with evaluative signals rather than embedding constraints
within its representational structure.

A proof-carrying world model offers a complementary approach. By requiring that
high-stakes action types carry typed witnesses of their preconditions, the
planning procedure is structurally prevented from producing plans that violate
stated constraints---not because a guard intercepts the output, but because the
output type cannot be constructed without the required evidence. The difference
is analogous to the difference between a runtime bounds check and a statically
verified array access: the latter eliminates a class of error by construction,
before execution begins.

We envision a planning module whose action type is parameterized by an evidence
profile $\epsilon$: an action of type $\mathsf{Action}(\epsilon)$ may only rely
on assumptions justified by $\epsilon$. Confidence escalation then becomes a
type error.

The JEPA connection matters here as well. If a predictive model is trained to
infer withheld structure from partial context, the question naturally follows:
what kinds of context are sufficient for which kinds of action? Dependent types
offer a way to state that question sharply. A latent prediction can be treated
not as a free-floating hint, but as an object whose admissible uses are tied to
the evidence from which it was formed.

\subsection{Multi-Agent Trust and Fragment Provenance}

In multi-agent settings, observations arrive from diverse sources with varying
degrees of trust. A dependently typed world model can represent this by
annotating observation histories with \emph{provenance types}: a history of
type $\mathsf{Hist}(\tau)$ was collected under trust level $\tau$, and the type
system can enforce that a planner operating under trust level $\tau'$ may only
act on histories whose provenance satisfies $\tau' \preceq \tau$ for some trust
ordering $\preceq$.

This makes cross-provenance contamination detectable at the type level---a
property with direct practical significance whenever agents interact with
adversarially constructed inputs or aggregated data of heterogeneous origin.

%% ─────────────────────────────────────────────────────────────────────────────
\section{Related Work}
\label{sec:related}
%% ─────────────────────────────────────────────────────────────────────────────

\paragraph{World models.}
Ha and Schmidhuber's World Models \citep{ha2018world} helped establish the
``dream-then-act'' paradigm for visual planning. The DreamerV1--V3 series
\citep{hafner2019dream, hafner2021mastering, hafner2023mastering} demonstrated
that continuous control and Atari tasks can be mastered through pure
imagination. MuZero \citep{schrittwieser2020mastering} showed that a latent
model without ground-truth reconstruction suffices for search-based planning.
EfficientZero \citep{ye2021efficientzero} extended this to extreme sample
efficiency. None of these architectures address typed invariants on the model
interface.

\paragraph{JEPA and predictive representation learning.}
LeCun's program for autonomous machine intelligence \citep{lecun2022path}
argues that useful intelligence should be grounded in predictive world models
that operate in latent space rather than through exhaustive reconstruction.
I-JEPA \citep{assran2023ijepa} makes this idea concrete by learning to predict
target embeddings from context embeddings. Our use of typed predictive
contracts is meant to be complementary: JEPA supplies a modern account of how
predictive latent structure may be learned, while Rocq supplies a way to state,
with precision, what such predictions are allowed to authorize in planning.

\paragraph{Formal methods and learning.}
Seshia et al.\ \citep{seshia2018formal} survey formal specification and
verification techniques applicable to learning-enabled systems, emphasizing
runtime monitors and assume-guarantee contracts. Neural network verification
\citep{katz2019marabou, liu2021algorithms} focuses on robustness of fixed
networks. CompCert \citep{leroy2009compcert} and seL4 \citep{klein2009sel4}
demonstrate that large software systems can be given machine-checked
correctness proofs; our work proposes that learned interfaces can participate
in such developments.

\paragraph{Alignment and specification.}
Russell \citep{russell2019human} frames alignment as a problem of
misspecified objectives. Christiano et al.\ \citep{christiano2017rlhf} and
Bai et al.\ \citep{bai2022training, bai2022constitutional} address alignment
through human feedback and constitutional constraints. Gabriel
\citep{gabriel2020alignment} provides a conceptual taxonomy. Our proposal is
orthogonal: rather than specifying objectives more accurately, we propose to
specify the operational interface of the model more precisely.

%% ─────────────────────────────────────────────────────────────────────────────
\section{Conclusion}
\label{sec:conclusion}
%% ─────────────────────────────────────────────────────────────────────────────

We have argued that world models become easier to reason about when they are
understood as accountable constructions rather than opaque latent containers.
JEPA reminds us that useful predictive structure can be learned from partial
views. Dependent type theory adds a discipline for saying what those predictions
license: when a state is admissible, when evidence is sufficient, and when a
plan should abstain.

The proposal is not that proofs replace learning. A well-typed world model may
still be empirically wrong about the world. The proposal is narrower and, for
that reason, practical: that learned internal models can carry clearer
commitments about admissibility, provenance, and use in multi-agent systems. If
a system is going to imagine on our behalf, it is reasonable to ask that its
imagination come with explicit terms.

%% ─────────────────────────────────────────────────────────────────────────────
\bibliographystyle{plainnat}
\bibliography{world_model_paper}
%% ─────────────────────────────────────────────────────────────────────────────

\end{document}

